\section{Introduction}
\label{sec:introduction}

\begin{itemize}

  % ----- Background -----
  \item \textbf{Background:}
  Panoramic indoor layout annotation underpins 3D reconstruction, robot navigation, and AR/VR.
  At scale, pipelines rely on \emph{semi-automatic crowd annotation}---a model proposes an initial
  layout, workers edit it---which introduces several compounding QA challenges.
  \begin{itemize}
    \item \textbf{Budget constraint:} exhaustive expert re-verification is infeasible; only a fraction of submissions can be reviewed.
    \item \textbf{Heterogeneity:} scene difficulty (e.g., occlusion, low-texture surfaces), model failure modes, and out-of-scope cases (e.g., split-level or multi-plane scenes) are not uniformly distributed; worker capabilities vary.
    \item \textbf{Instrumentation gap:} standard annotation platforms do not natively record true \emph{active engagement time} or structured per-task diagnostic signals, making efficiency and difficulty claims hard to verify.
  \end{itemize}

  % ----- Motivation -----
  \item \textbf{Motivation:}
  Without pre-specified task distribution rules and locked thresholds, practitioners often default to
  \emph{retrospective} filtering---deciding what to reject or re-annotate \emph{after} seeing results.
  This post-hoc approach relies on undocumented threshold selection and risks
  \emph{metric inflation via silent pruning}: hard cases (which cluster under distribution shift)
  are quietly removed, making reported metrics reflect an easy subset rather than the full distribution.
  \begin{itemize}
    \item \textbf{Need 1a:} Evaluate whether semi-automatic initialization genuinely reduces active annotation time and improves consistency under matched, controlled conditions.
    \item \textbf{Need 1b:} Assess whether pre-screening and risk-aware task distribution can sustain reliability under non-i.i.d./OOD scenes while improving budget efficiency.
    \item \textbf{Need 2:} Ensure gate failures and missing-field events are disclosed as explicit audit tiers rather than silently absorbed, so that reported metrics reflect the full data distribution.
  \end{itemize}

  % ----- Our Approach -----
  \item \textbf{Our Approach:}
  We extend Label Studio with a custom userscript that captures precise active engagement time and
  collects per-task annotator judgments on \texttt{scope}, \texttt{difficulty}, and \texttt{model\_issue}.
  On top of this instrumented platform we propose a pre-registered, staged protocol
  (Pilot $\rightarrow$ PreScreen $\rightarrow$ Main) covering \emph{Pre-screening, Calibration, Task Distribution, and Audit}. Pilot locks the protocol skeleton, PreScreen locks admission-related values, and the first Calibration round locks routing parameters that genuinely depend on realized coverage.
  \begin{itemize}
    \item \textbf{Audit-first workflow:} pre-annotation routing relies on cold-start risk proxies such as $d_t$ and $g_t$, while meta-labels primarily serve as post-submission audit and scene-stratification signals; outcomes are reported under transparent disclosure tiers (Total / In-scope / Model-usable, i.e., T/I/M).
    \item \textbf{Risk-aware task distribution:} effort is allocated adaptively for high-risk/OOD tasks via the distribution strategy $\pi$; budget is controlled via sequential redundancy and stopping rules.
    \item \textbf{Comparable evaluation:} task distribution strategies are compared via offline replay on tasks with multiple real annotations, avoiding counterfactual label imputation.
  \end{itemize}

  % ----- Contributions -----
  \item \textbf{Contributions:}
  \begin{itemize}
    \item A customized annotation platform (Label Studio + userscript) that records verified active engagement time and structured per-task diagnostic meta-labels.
    \item An end-to-end, pre-registered QA protocol for semi-automatic panoramic layout annotation with explicit audit outputs.
    \item A transparent reporting scheme (T/I/M) that treats gate failures and missing fields as safety/data-availability audit items, preventing silent metric inflation.
    \item A risk-aware annotation task distribution framework targeting budget efficiency under non-i.i.d./OOD conditions, with strategy comparison via offline replay.
  \end{itemize}

\end{itemize}
